\section{МЕТОДИЧНЕ ЗАБЕЗПЕЧЕННЯ НАВЧАННЯ РОЗРОБКИ ІМЕРСИВНИХ ОСВІТНІХ РЕСУРСІВ}

У розділі розкрито організаційно-педагогічні засади навчання розробки імерсивних освітніх ресурсів майбутніх учителів інформатики у контексті цифрової трансформації освіти. Проаналізовано наукові підходи та принципи, які визначають методологічні засади підготовки здобувачів до створення VR/AR-рішень і моделювання інтерактивних навчальних середовищ. Запропоновано модель формування вмінь створювати імерсивні освітні ресурси у процесі фахової підготовки. Окреслено педагогічні умови, що забезпечують ефективність цього процесу.

\subsection{Дидактичні підходи та принципи навчання розробки імерсивних освітніх ресурсів майбутніх учителів інформатики}\label{subsec:2.1}

Розвиток імерсивних технологій відкриває нові можливості для цифрової трансформації педагогічної освіти, зокрема у сфері підготовки майбутніх учителів інформатики. Сучасна педагогічна парадигма орієнтує освітній процес на формування не лише цифрових, а й проєктно-технологічних компетентностей, що дозволяють створювати інноваційні освітні продукти у форматі VR/AR/MR. Осмислення методологічних основ навчання студентів розробки імерсивних освітніх ресурсів потребує визначення комплексу наукових підходів, здатних забезпечити наукову цілісність і практичну ефективність підготовки.

У своєму розумінні дидактичних підходів та принципів навчання майбутніх учителів виходимо з сучасного розуміння понять «дидактичний підхід» та «принципи навчання». У науково-педагогічному дискурсі дослідники пропонують широкий перелік дидактичних підходів, які відображають різні методологічні та методичні орієнтири професійної підготовки. У сучасній педагогічній літературі акцентується значущість знаннєво орієнтованого, особистісно зорієнтованого, діяльнісно-розвивального та компетентнісного підходів~\cite{Kuzminskyi2012Pedagogy}, аксіологічного та культурологічного підходів як підґрунтя розвитку професійної ідентичності~\cite{Kremen2016Axiological}, інноваційного підходу, що передбачає впровадження нових форм і методів навчання~\cite{Dychkivska2012Innovative}, а також системного, прогностичного та компетентнісного підходів, що забезпечують цілісність освітнього процесу~\cite{Sysoieva2015Pedagogy}.

Вибір дидактичних підходів у нашому дослідженні спирається на теоретичні результати першого розділу та має дедуктивний характер. Когнітивно-афективна модель імерсивного навчання (CAMIL)~\cite{Makransky2021CAMIL} визначає шість медіативних факторів ефективності імерсивних середовищ~-- присутність, агентність, втілення, когнітивне навантаження, мотивацію та саморегуляцію,~-- кожен із яких потребує відповідного методологічного забезпечення. Водночас виявлений у першому розділі дефіцит технолого-педагогічних знань (TPK-компонент моделі TPACK)~\cite{Gurevych2015Integration} зумовлює потребу в підходах, що інтегрують технологічну й педагогічну підготовку. Нарешті, п'ять сутнісних характеристик ІОР, обґрунтованих у підрозділі~1.2 (інтерактивність, імерсивність, контекстуальність, візуалізація, адаптивність), задають критерії повноти обраної сукупності підходів. Систематичний огляд~\cite{Mallek2024SynthesizingVR} підтверджує, що ефективність імерсивних навчальних середовищ зростає за умови узгодження педагогічного проєктування з когнітивними теоріями навчання.

З огляду на ці теоретичні засади, а також з урахуванням специфіки предмета й об'єкта нашого дослідження, було виокремлено одинадцять концептуально важливих підходів~-- аксіологічний, діяльнісний, інноваційний, інтегративний, когнітивний, компетентнісний, конструктивістський, контекстний, рефлексивний, синергетичний і системний,~-- сукупність яких забезпечує покриття усіх шести факторів CAMIL, усунення TPK-дефіциту та відповідність п'яти характеристикам ІОР.

Розглянемо окреслені дидактичні підходи більш детально, оскільки вони формують теоретичне підґрунтя ефективної підготовки майбутніх учителів інформатики до розробки ІОР.

\emph{Аксіологічний підхід} «полягає в гуманістичній спрямованості освітнього процесу на розвиток особистості як мету, результат і головний критерій її ефективності»~\cite{Kremen2016Axiological}. Звертаємо увагу і на твердження про те, що «аксіологічний підхід визнає особистісно значущі та життєво важливі потреби людини як найвищу цінність»~\cite[c.~49]{Kuzminskyi2012Pedagogy}, орієнтуючи освітній процес на саморозвиток суб'єктів і запобігаючи формуванню інтуїтивних уявлень про цінності.

У підготовці майбутніх учителів інформатики до розробки ІОР аксіологічний підхід підкреслює ціннісний вимір професійної діяльності, зокрема відповідальність за створення цифрового освітнього контенту, що відповідає етичним нормам (коректність подання інформації, уникнення маніпулятивних елементів, безпечна взаємодія користувачів із цифровим середовищем), принципам гуманізму, поваги до особистості, відповідального використання технологій та академічної доброчесності. У контексті моделі CAMIL підхід безпосередньо адресує фактор мотивації, оскільки ціннісне осмислення власної діяльності сприяє формуванню внутрішніх мотивів створення якісного імерсивного контенту~\cite{Lee2022VRMaking}. Важливим стає усвідомлення здобувачами наслідків застосування створених ними VR/AR-ресурсів у реальному освітньому середовищі.

\emph{Діяльнісний підхід} передусім пов'язаний із організацією діяльності, яка забезпечує розвиток і становлення особистості майбутнього вчителя інформатики~\cite{Dychkivska2012Innovative} та визначає орієнтири його практичної підготовки до розробки ІОР. Підхід акцентує увагу на поступовому формуванні навичок через виконання практикоорієнтованих завдань, у яких здобувач виступає активним суб'єктом створення імерсивного контенту. Замість пасивного сприйняття інформації, здобувачі активно залучаються до процесу розробки VR/AR-рішень, експериментують з різними інструментами та технологіями, вирішують професійні проблеми та приймають рішення, що сприяє грунтовному засвоєнню матеріалу та розвитку практичних навичок. Емпіричне дослідження~\cite{Kee2023Experiential} підтверджує, що діяльнісна організація навчання в імерсивних середовищах за циклом Колба (конкретний досвід~-- рефлексивне спостереження~-- абстрактна концептуалізація~-- активне експериментування) суттєво підвищує якість засвоєння проєктних умінь; у контексті CAMIL такий підхід забезпечує фактор агентності.

\emph{Інноваційний підхід} забезпечує відкритість системи підготовки до технологічних змін, інтеграцію новітніх інструментів і платформ у процес навчання. У наукових джерелах інноваційний підхід трактується як «єдиний змістовий комплекс принципів і методів, що використовуються в освітньому процесі та спрямовані на зміну вимог до якості освіти»~\cite[c.~33]{Sysoieva2015Pedagogy}. Підхід зумовлює необхідність гнучкої й адаптивної організації підготовки, здатної оперативно інтегрувати новітні інструменти, платформи та технологічні рішення. У підготовці майбутніх учителів інформатики інноваційний підхід орієнтує на формування готовності працювати з новими цифровими засобами, оскільки передбачає «впровадження в освітній процес інформаційно-комп'ютерних технологій, інноваційних форм і методів навчання»~\cite[c.~34]{Sysoieva2015Pedagogy}, що створює умови для практичного опанування сучасних VR/AR-платформ, конструкторів, програмних середовищ та їх використання під час розробки імерсивних освітніх продуктів. Як засвідчує систематичний огляд~\cite{Crogman2025ExperientialXR}, XR-технології трансформують освітню парадигму саме завдяки інноваційному поєднанню досвідного навчання з імерсивними середовищами.

\emph{Інтегративний підхід} забезпечує цілісність підготовки майбутніх учителів інформатики до розробки ІОР, дозволяє поєднати технічну, дидактичну та психологічну компоненти підготовки в єдину систему. Створення ІОР вимагає не лише технічних знань, але й розуміння педагогічних принципів та психологічних особливостей учнів. Інтегративний підхід забезпечує гармонійне поєднання цих компонентів, коли майбутній учитель одночасно опановує технологічні інструменти, методичні принципи та особливості сприймання інформації учнями, що дозволяє створювати ефективні та цікаві навчальні матеріали. Цю позицію підтримує концепція «переплетеного еклектизму» (entangled eclecticism)~\cite{Schmidt2024EntangledLXD}, згідно з якою ефективний дизайн навчального досвіду вимагає одночасного врахування соціокультурних, технологічних і педагогічних вимірів як взаємопов'язаної екології, що безпосередньо кореспондує з TPK-компонентом моделі TPACK.

У контексті професійної підготовки \emph{когнітивний підхід} передбачає орієнтацію на активну розумову діяльність здобувачів освіти, що узгоджується з концепцією Д.~Шона про рефлексію в дії (reflection-in-action) та рефлексію над дією (reflection-on-action), яка акцентує здатність фахівця до осмислення власних когнітивних процесів безпосередньо під час виконання професійних завдань і після їх завершення~\cite{ReflexiveApproach2018}. Когнітивний підхід в аспекті нашої теми робить акцент на процесах сприймання, осмислення й опрацювання інформації, що реалізується через використання візуалізації, моделювання та симуляцій як засобів поглиблення розуміння навчального матеріалу. Систематичний огляд~\cite{Mallek2024SynthesizingVR} засвідчує, що результати навчання у VR-середовищах суттєво покращуються за умови узгодження педагогічного дизайну з теоріями когнітивного навантаження; метааналіз впливу доповненої реальності на когнітивне навантаження~\cite{Buchner2021CogLoad} підтверджує необхідність свідомого управління зовнішнім когнітивним навантаженням при проєктуванні AR-ресурсів. Таким чином, у контексті моделі CAMIL когнітивний підхід безпосередньо адресує фактор когнітивного навантаження.

\emph{Компетентнісний підхід} забезпечує спрямованість підготовки на формування інтегрованих професійних умінь~-- проєктування, програмування, педагогічного дизайну та цифрової креативності. За таких умов здобувачі мають не лише опановувати теоретичні знання, а й набувати здатності до їх практичного застосування у процесі створення реальних імерсивних освітніх продуктів, що відповідає розумінню компетентності як \textbf{«}спроможності кваліфіковано здійснювати діяльність, виконувати завдання або роботу, поєднуючи знання, уміння, навички та ставлення»~\cite[c.~240]{Dychkivska2012Innovative}. У межах компетентнісного підходу особливого значення набуває оволодіння навичками проєктування сценаріїв, інтерактивного програмування, розробки візуального дизайну та креативного вирішення педагогічних завдань, оскільки професійна компетентність трактується як «інтегративна якість особистості, що виявляється у здатності ефективно діяти на основі практичного досвіду, знань і вмінь»~\cite[c.~96]{Dychkivska2012Innovative}. Дослідження~\cite{BeldaMedina2022TPACK} емпірично підтверджує, що формування TPACK-компетентності через практичне створення AR-ресурсів суттєво підвищує цифрову компетентність і самоефективність майбутніх учителів.

\emph{Конструктивістський підхід} передбачає активне залучення здобувачів до побудови власного знання через створення цифрових освітніх продуктів і взаємодію у віртуальному середовищі. Ми зважаємо на думку В.~Соловйова із співавторами, про те, що «конструктивізм виходить з того, що навчання~-- це активний процес, у ході якого суб'єкти активно конструюють знання на основі власного досвіду. Ідеї конструктивізму виражені й у теорії діяльності, згідно з якою діяльність і дії є основою психічного розвитку»~\cite{Teplytskyi2015Modelling}. Тобто здобувачі не просто отримують готові знання, а й самостійно визначають логіку, структуру та функціональне призначення імерсивних ресурсів, активно конструюючи їх. Систематичний огляд~\cite{Mallek2024SynthesizingVR} визначає конструктивізм як найбільш ефективне педагогічне підґрунтя для імерсивного навчання, а дослідження VR-Making~\cite{Lee2022VRMaking} демонструє, що конструкціоністська діяльність зі створення VR-середовищ розвиває готовність учителів до технологічної інтеграції, що безпосередньо адресує фактор втілення у моделі CAMIL.

\emph{Контекстний підхід} спрямовує навчання у професійне поле, забезпечуючи його максимальну наближеність до реальних умов педагогічної діяльності та відтворення ситуацій шкільної практики, у межах яких ІОР набуває конкретного функціонального призначення. Така орієнтація відповідає сутності контекстного навчання, яке «проектує освітній процес у вищому навчальному закладі як максимально наближений до майбутньої професійної діяльності»~\cite{Spirin2020Ref}. Це також сприяє розумінню професійної підготовки вчителя як процесу формування здатності до усвідомленого виконання педагогічних дій у реальних освітніх умовах, оскільки підготовка майбутнього вчителя має забезпечувати якісне виконання педагогічних дій через формування загальних і спеціальних компетентностей~\cite{Dychkivska2012Innovative}.

У процесі підготовки майбутні вчителі інформатики залучаються до моделювання педагогічних ситуацій, що дає змогу усвідомити, як створені VR/AR-рішення використовуватимуться в реальному навчальному середовищі, з якими освітніми завданнями вони співвідносяться та які педагогічні ефекти мають забезпечувати, оскільки «контекстним навчанням є навчання, в якому динамічно моделюється предметний та соціальний зміст професійної діяльності»~\cite{Spirin2020Ref}. Емпіричні дані підтверджують цю позицію: дослідження~\cite{Li2024ClassMasterIVR} засвідчує, що повністю імерсивне віртуальне середовище суттєво підвищує професійну компетентність майбутніх учителів саме через автентичність педагогічної практики, а огляд~\cite{Paulsen2024CollaborativeIVR} визначає імерсивну віртуальну реальність як «простір практики та рефлексії», що забезпечує фактор присутності у моделі CAMIL. Загалом контекстний підхід сприяє формуванню професійної готовності до педагогічно виваженого використання VR/AR-технологій у шкільному навчанні~\cite{Kuzminskyi2012Pedagogy}.

\emph{Рефлексивний підхід} спрямований на формування здатності усвідомлювати власну діяльність у процесі розробки ІОР, критично оцінювати її результати та визначати напрями подальшого професійного вдосконалення. У науково-педагогічному дискурсі рефлексія розглядається як осмислення фахівцем власних дій та їх підстав, що, за Д.~Шоном, реалізується через рефлексію в дії та рефлексію над дією~\cite{ReflexiveApproach2018}. Застосування рефлексивного підходу сприяє усвідомленню здобувачем логіки та результативності власних дій на всіх етапах створення імерсивних продуктів, формує здатність аргументовано оцінювати обрані технічні й дидактичні рішення та співвідносити їх із поставленими освітніми цілями. Систематичний аналіз сильних і слабких сторін розроблених VR/AR-рішень, типових проєктних помилок і обмежень стимулює розвиток професійно значущих умінь самокорекції та відповідального ставлення до якості освітнього продукту. Водночас рефлексивний підхід створює підґрунтя для планування подальших кроків з удосконалення імерсивних продуктів, оптимізації індивідуального проєктного процесу та вибудовування навчально-професійної траєкторії майбутнього вчителя інформатики. Як підкреслює С.~Гончаренко, систематичне осмислення результатів педагогічної діяльності є необхідною умовою професійного зростання та вдосконалення педагогічної практики~\cite{Goncharenko2011Pedagogy}. Метааналіз~\cite{Han2025VRTeacherEd} засвідчує статистично значущий вплив VR на рефлексивну практику майбутніх учителів (\textit{g}=0,524), а дослідження~\cite{Kee2023Experiential} підтверджує, що рефлексивне спостереження та абстрактна концептуалізація (етапи RO та AC за Колбом) є критичними для перетворення імерсивного досвіду на стійкі професійні уміння, що забезпечує фактор саморегуляції у моделі CAMIL. Відтак рефлексивний підхід у підготовці майбутніх учителів інформатики до розробки ІОР виступає методологічним підґрунтям формування здатності до самостійного аналізу, самоконтролю та прогнозування власного професійного розвитку.

У наукових джерелах підкреслюється: синергетика міждисциплінарний напрям, який «розглядає освітні системи як відкриті, здатні до самокоригування та пошуку альтернативних шляхів розвитку»~\cite[c.~49]{Kuzminskyi2012Pedagogy}, що зумовлює переосмислення пізнавального процесу в педагогічній освіті. \emph{Синергетичний підхід} розкриває процес навчання як самоорганізовану систему, що з позиції відкритості, співтворчості та орієнтації на саморозвиток сприяє розв'язанню протиріч, що виникають на різних рівнях освіти~\cite{Kuzminskyi2012Pedagogy}, у якій нові знання виникають завдяки взаємодії здобувача, технології та освітнього середовища.

У межах імерсивного навчання синергетичні ідеї набувають особливої значущості, адже нові знання формуються в результаті взаємодії здобувача освіти, цифрових технологій та освітнього середовища, у якому технологія постає активним учасником освітньої взаємодії, стимулюючи виникнення індивідуальних освітніх траєкторій. Цю позицію підтримує системна теорія дизайну навчального досвіду~\cite{Schmidt2024EntangledLXD}, яка описує освітнє середовище як «взаємопов'язану екологію» соціокультурних, технологічних і педагогічних вимірів, де емерджентні ефекти виникають саме завдяки їхній взаємодії. Відповідно синергетичний підхід дозволяє розглядати навчання розробки VR/AR-рішень як відкриту самоорганізовану систему, що забезпечує фактор присутності у моделі CAMIL через поєднання технологічних, когнітивних і дидактичних компонентів.

Науковці практично одностайні в тому, що системний підхід передбачає розгляд об'єкта як «цілісної системи з визначенням її елементів та упорядкуванням зв'язків між ними»~\cite[c.~46]{Kuzminskyi2012Pedagogy}, водночас акцентуючи на необхідності виокремлення системотвірних взаємодій і врахування впливу надсистеми та внутрішньої взаємодії підсистем. Така позиція дає підстави розглядати системний підхід як інструмент виявлення причинно-наслідкових залежностей і наукового узагальнення~\cite{Goncharenko2011Pedagogy}.

У межах дослідження системний підхід дає змогу розглядати процес розробки ІОР як цілісну педагогічну систему зі стійкими структурними й функціональними зв'язками, у якій усі компоненти перебувають у взаємозумовленій взаємодії. Його застосування уможливлює побудову моделі навчання, у межах якої мета, зміст, методи, засоби та результат поєднані в логічно впорядковану структуру, що забезпечує узгодженість освітніх впливів, оптимізацію навчального процесу та підвищення його ефективності й результативності. У такій логіці системний підхід виступає одним із провідних методологічних засобів організації підготовки майбутніх учителів інформатики до розробки ІОР, оскільки сприяє комплексному опануванню технологій створення VR/AR-матеріалів як цілісної системи, що відображає взаємопов'язаність програмних, дидактичних і дизайнерських компонентів професійної діяльності. Системний підхід також орієнтує на свідомий добір оптимальних способів розв'язання навчально-пізнавальних і проєктних завдань, визначення адекватних теоретичних і практичних засобів підтримки розробки ІОР. Як підтверджує дослідження~\cite{Schmidt2024EntangledLXD}, соціокультурний, технологічний і педагогічний виміри освітнього середовища мають розглядатися як «зв'язна система» (cohesive system), у якій зміна одного компонента впливає на всі інші, що підсилює методичну вивіреність і педагогічну доцільність створюваних VR/AR-рішень та адресує TPK-інтеграцію у моделі TPACK.

Функціональне призначення використання тих чи тих підходів у розробці методики навчання розробки ІОР майбутніх учителів інформатики унаочнено на рис.~\ref{fig:2.1}.

\begin{figure}[!h]
\centering
%\resizebox{\textwidth}{!}{%
\begin{tikzpicture}[node distance=0.5cm and 0.6cm]
\small

% Title
\node[dia/header=15cm] (title) {ДИДАКТИЧНІ ПІДХОДИ};

% === LEFT COLUMN (3 approaches) ===
\node[dia/lrbox=4.5cm, below=2.0cm of title.south, xshift=-5.3cm] (a1) {%
  \textbf{Аксіологічний}\\[2pt]
  {\small Забезпечує етичну та гуманістичну розробку контенту}%
};
\node[dia/lrbox=4.5cm, below=0.5cm of a1] (a2) {%
  \textbf{Синергетичний}\\[2pt]
  {\small Сприяє новим знанням через взаємодію між здобувачами, технологіями та середовищем}%
};
\node[dia/lrbox=4.5cm, below=0.5cm of a2] (a3) {%
  \textbf{Компетентнісний}\\[2pt]
  {\small Фокусується на розвитку інтегрованих професійних навичок, необхідних для створення ІОР}%
};

% === CENTER COLUMN (5 approaches) ===
\node[dia/lrbox=4.5cm, below=2.0cm of title.south] (b1) {%
  \textbf{Конструктивістський}\\[2pt]
  {\small Залучає здобувачів до побудови власного знання}%
};
\node[dia/lrbox=4.5cm, below=0.5cm of b1] (b2) {%
  \textbf{Інноваційний}\\[2pt]
  {\small Інтегрує новітні технології та інструменти}%
};
\node[dia/lrbox=4.5cm, below=0.5cm of b2] (b3) {%
  \textbf{Когнітивний}\\[2pt]
  {\small Враховує когнітивні процеси при проєктуванні навчального досвіду}%
};
\node[dia/lrbox=4.5cm, below=0.5cm of b3] (b4) {%
  \textbf{Контекстний}\\[2pt]
  {\small Адаптує ресурси до реальних освітніх ситуацій}%
};
\node[dia/lrbox=4.5cm, below=0.5cm of b4] (b5) {%
  \textbf{Рефлексивний}\\[2pt]
  {\small Заохочує самооцінку та вдосконалення}%
};

% === RIGHT COLUMN (3 approaches) ===
\node[dia/lrbox=4.5cm, below=2.0cm of title.south, xshift=5.3cm] (c1) {%
  \textbf{Системний}\\[2pt]
  {\small Розглядає розробку як цілісну систему з взаємопов'язаними компонентами}%
};
\node[dia/lrbox=4.5cm, below=0.5cm of c1] (c2) {%
  \textbf{Інтегративний}\\[2pt]
  {\small Поєднує технічні, дидактичні та психологічні компоненти цілісного навчання}%
};
\node[dia/lrbox=4.5cm, below=0.5cm of c2] (c3) {%
  \textbf{Діяльнісний}\\[2pt]
  {\small Підкреслює практичне навчання через активну участь у розробці}%
};

% Lines from title to column tops
\draw[dia/line] (title.south) -- (a1.north);
\draw[dia/line] (title.south) -- (b1.north);
\draw[dia/line] (title.south) -- (c1.north);

\end{tikzpicture}%
%}

\caption{Дидактичні підходи до формування методики навчання розробки імерсивних освітніх ресурсів}
\label{fig:2.1}
\end{figure}

Сукупність поданих \emph{дидактичних підходів} утворює методологічну основу підготовки майбутніх учителів інформатики до розробки імерсивних освітніх ресурсів. Застосування сукупності підходів формує новий тип педагога~-- інженера освітніх технологій, здатного інтегрувати технологічні інновації з педагогічною доцільністю й гуманістичними цінностями освіти.

Для обґрунтування повноти обраної сукупності підходів побудовано матрицю їх відповідності п'яти сутнісним характеристикам ІОР (підрозділ~1.2) та шести медіативним факторам моделі CAMIL (рис.~\ref{fig:approaches_matrix}). Матриця демонструє, що кожна характеристика ІОР забезпечується щонайменше трьома підходами, а кожен фактор CAMIL адресується щонайменше двома підходами, що підтверджує теоретичну обґрунтованість і повноту запропонованої сукупності.

\begin{figure}[!h]
\centering
\resizebox{\textwidth}{!}{%
\begin{tikzpicture}
\footnotesize

% Column headers (IER characteristics)
\def\colW{2.4cm}
\def\rowH{0.7cm}
\def\hdrH{1.2cm}
\def\startX{4.8cm}

% Header row — IER characteristics
\node[rectangle, draw, fill=black!10, text width=\colW, minimum height=\hdrH,
      text centered, font=\bfseries\scriptsize] at (\startX, 0)
  {Інтерак\-тивність};
\node[rectangle, draw, fill=black!10, text width=\colW, minimum height=\hdrH,
      text centered, font=\bfseries\scriptsize] at (\startX + 2.6cm, 0)
  {Імерсив\-ність};
\node[rectangle, draw, fill=black!10, text width=\colW, minimum height=\hdrH,
      text centered, font=\bfseries\scriptsize] at (\startX + 5.2cm, 0)
  {Контексту\-альність};
\node[rectangle, draw, fill=black!10, text width=\colW, minimum height=\hdrH,
      text centered, font=\bfseries\scriptsize] at (\startX + 7.8cm, 0)
  {Візуалі\-зація};
\node[rectangle, draw, fill=black!10, text width=\colW, minimum height=\hdrH,
      text centered, font=\bfseries\scriptsize] at (\startX + 10.4cm, 0)
  {Адаптив\-ність};

% Extra column: CAMIL factors
\node[rectangle, draw, fill=black!15, text width=3.2cm, minimum height=\hdrH,
      text centered, font=\bfseries\scriptsize] at (\startX + 13.4cm, 0)
  {Фактори\\CAMIL};

% Row label column header
\node[rectangle, draw, fill=black!10, text width=2.8cm, minimum height=\hdrH,
      text centered, font=\bfseries\scriptsize] at (2.0cm, 0)
  {Дидактичний\\підхід};

% --- Data rows ---
% Approach names (left column)
\foreach \i/\name in {%
  1/Аксіологічний,
  2/Діяльнісний,
  3/Інноваційний,
  4/Інтегративний,
  5/Когнітивний,
  6/Компетентнісний,
  7/Конструктивістський,
  8/Контекстний,
  9/Рефлексивний,
  10/Синергетичний,
  11/Системний%
}{%
  \node[rectangle, draw, text width=2.8cm, minimum height=\rowH,
        inner sep=3pt, font=\scriptsize] at (2.0cm, -\i*0.8cm)
    {\name};
}

% Dot markers: approach × characteristic
% Format: \fill (x, y) circle (3pt);
% Col positions: Inter=\startX, Immers=+2.6, Context=+5.2, Visual=+7.8, Adapt=+10.4

% Аксіологічний: Контекстуальність
\fill (\startX + 5.2cm, -0.8cm) circle (3pt);

% Діяльнісний: Інтерактивність, Контекстуальність
\fill (\startX, -1.6cm) circle (3pt);
\fill (\startX + 5.2cm, -1.6cm) circle (3pt);

% Інноваційний: Адаптивність, Візуалізація
\fill (\startX + 7.8cm, -2.4cm) circle (3pt);
\fill (\startX + 10.4cm, -2.4cm) circle (3pt);

% Інтегративний: Інтерактивність, Імерсивність, Візуалізація
\fill (\startX, -3.2cm) circle (3pt);
\fill (\startX + 2.6cm, -3.2cm) circle (3pt);
\fill (\startX + 7.8cm, -3.2cm) circle (3pt);

% Когнітивний: Візуалізація, Імерсивність
\fill (\startX + 7.8cm, -4.0cm) circle (3pt);
\fill (\startX + 2.6cm, -4.0cm) circle (3pt);

% Компетентнісний: Інтерактивність, Контекстуальність, Адаптивність
\fill (\startX, -4.8cm) circle (3pt);
\fill (\startX + 5.2cm, -4.8cm) circle (3pt);
\fill (\startX + 10.4cm, -4.8cm) circle (3pt);

% Конструктивістський: Інтерактивність, Імерсивність
\fill (\startX, -5.6cm) circle (3pt);
\fill (\startX + 2.6cm, -5.6cm) circle (3pt);

% Контекстний: Імерсивність, Контекстуальність
\fill (\startX + 2.6cm, -6.4cm) circle (3pt);
\fill (\startX + 5.2cm, -6.4cm) circle (3pt);

% Рефлексивний: Адаптивність, Контекстуальність
\fill (\startX + 10.4cm, -7.2cm) circle (3pt);
\fill (\startX + 5.2cm, -7.2cm) circle (3pt);

% Синергетичний: Імерсивність, Адаптивність
\fill (\startX + 2.6cm, -8.0cm) circle (3pt);
\fill (\startX + 10.4cm, -8.0cm) circle (3pt);

% Системний: Інтерактивність, Візуалізація, Контекстуальність
\fill (\startX, -8.8cm) circle (3pt);
\fill (\startX + 7.8cm, -8.8cm) circle (3pt);
\fill (\startX + 5.2cm, -8.8cm) circle (3pt);

% CAMIL factor labels (right column)
\foreach \i/\camil in {%
  1/{мотивація},
  2/{агентність},
  3/{агентність},
  4/{когн. навант.},
  5/{когн. навант.},
  6/{саморегуляція},
  7/{втілення},
  8/{присутність},
  9/{саморегуляція},
  10/{присутність},
  11/{когн. навант.}%
}{%
  \node[rectangle, draw, text width=3.2cm, minimum height=\rowH,
        inner sep=3pt, font=\scriptsize, text centered] at (\startX + 13.4cm, -\i*0.8cm)
    {\camil};
}

% Grid lines (horizontal)
\foreach \i in {1,...,11}{%
  \draw[gray!40] (0cm, -\i*0.8cm + 0.4cm) -- (\startX + 15.2cm, -\i*0.8cm + 0.4cm);
  \draw[gray!40] (0cm, -\i*0.8cm - 0.4cm) -- (\startX + 15.2cm, -\i*0.8cm - 0.4cm);
}

% Grid lines (vertical)
\foreach \dx in {0cm, 2.6cm, 5.2cm, 7.8cm, 10.4cm}{%
  \draw[gray!40] (\startX + \dx - 1.3cm, 0.7cm) -- (\startX + \dx - 1.3cm, -9.2cm);
  \draw[gray!40] (\startX + \dx + 1.3cm, 0.7cm) -- (\startX + \dx + 1.3cm, -9.2cm);
}

\end{tikzpicture}%
}

\caption{Матриця відповідності дидактичних підходів характеристикам ІОР та факторам CAMIL}
\label{fig:approaches_matrix}
\end{figure}

Ми поділяємо позицію вчених, які стверджують, що реалізація будь-якого дидактичного підходу передбачає добір відповідних принципів навчання~\cite{Dychkivska2012Innovative}. У цьому контексті зосередимо увагу на принципах, які корелюють із визначеними вище дидактичними підходами та формують методологічне підґрунтя навчання розробки імерсивних освітніх ресурсів майбутніх учителів інформатики.

Аналізуючи структуру та функції дидактичних принципів у професійній підготовці фахівців, зокрема педагогів, науковці акцентують на їхньому значенні для забезпечення цілісності, системності й ефективності освітнього процесу. Представлені в сучасній літературі підходи до трактування дидактичних принципів дозволяють окреслити методологічне поле нашого дослідження, а також визначити ті з них, що найбільш повно відповідають специфіці підготовки здобувачів до розробки імерсивних освітніх ресурсів.

У сучасному науково-педагогічному дискурсі дидактичні принципи інтерпретуються як методологічні \emph{засади}, що визначають зміст і логіку професійної підготовки. І.~Дичківська виокремлює \emph{загальнодидактичні принципи} (науковість, наочність, послідовність та систематичність, зв'язок теорії з практикою, самостійності та активності), а також \emph{специфічні для фахової підготовки принципи} (проблемності та професійної мобільності), реалізація яких спрямована на формування професійних умінь через ускладнення навчальних завдань, активізацію мислення та орієнтацію на реальні професійні ситуації~\cite{Dychkivska2012Innovative}. У близькому методологічному ключі науковці розглядають принципи професійної підготовки майбутніх педагогів засобами ІКТ, акцентуючи увагу на інтерактивності, адаптивності, гнучкості, мультимедійності та співпраці~\cite{Semerikov2023Motivational}.

Позицію \emph{системності} й \emph{наукової цілісності} у трактуванні дидактичних принципів обґрунтовує С.~Гончаренко~\cite{Goncharenko2011Pedagogy}, який розглядає принципи як структурні елементи моделі професійної підготовки, здатні забезпечити узгодженість між цілями, змістом, методами й формами навчання. У межах системного підходу виокремлюються принципи \emph{науковості, системності, послідовності, безперервності, автономності навчання та зв'язку з практикою.} Широкий спектр дидактичних принципів навчання (\emph{систематичність та послідовність, науковість, професійне спрямування, наочність, посильність, міцність, свідомість, активність, виховне навчання, доступність, врахування індивідуальних особливостей та міжкультурної взаємодії)} висвітлюється також у працях І.~Дичківської~\cite{Dychkivska2012Innovative}.

Звернення до наукового доробку вчених дозволяє окреслити базовий перелік дидактичних принципів, які забезпечують методологічну цілісність професійної підготовки. Водночас специфіка навчання розробки ІОР майбутніх учителів інформатики потребує уточнення цього переліку з урахуванням технологічної, когнітивної та проєктної складових діяльності. Тому далі розглянемо групу принципів, що, на нашу думку, найбільш повно відповідають завданням підготовки здобувачів до створення VR/AR-ресурсів та забезпечують результативність цього процесу.

Визначення принципів навчання спирається на взаємозв'язок методологічних підходів, що покладений в основу формування цілісної системи методичних засад підготовки майбутніх учителів інформатики до розробки імерсивних освітніх ресурсів. Специфіка цього процесу зумовила виокремлення групи принципів \emph{інтерактивності, науковості, практичної спрямованості, креативності, варіативності, інтеграції, рефлексивності} та \emph{безперервності}, сукупність яких забезпечують якість та результативність процесу навчання.

\emph{Принцип інтерактивності} орієнтує на активну діяльність студента як творця цифрового освітнього контенту. Взаємодія з інструментами VR/AR-розробки відбувається з іншими учасниками навчального процесу. Вихідним положенням якого стає посилення мотиваційного потенціалу~\cite{Semerikov2023Motivational}. Дослідники зазначають, що саме інтерактивність імерсивних середовищ сприяє глибшому залученню студентів до професійної діяльності~\cite{Semerikov2022IERDesign}. С.~Литвинова разом із співавторами описують створення AR-об'єктів із мультимедійними шарами, що «оживляє» підручник через ігри, квести та симуляції~\cite[c.~56--57]{Lytvynova2023Monograph}.

\emph{Принцип науковості} забезпечує обґрунтованість розроблюваних освітніх ресурсів через опору на сучасні положення когнітивної психології, педагогічного дизайну та цифрової педагогіки. Наукова достовірність контенту виступає ключовою умовою якості імерсивних освітніх середовищ, зокрема у процесі моделювання складних явищ і процесів. С.~Литвинова та співавтори зазначають, що VR/AR-технології за дотримання змістових вимог забезпечують відтворення процесів із високим рівнем наукової точності, наголошуючи, що «контент має відповідати цілому набору вимог, у тому числі науковій достовірності»~\cite[c.~23--24; 26]{Lytvynova2023Monograph}. У методичному вимірі принцип науковості розглядається як узгодженість змісту навчання з актуальним станом розвитку науки та логікою наукового пізнання, що підвищує концептуальну точність і дидактичну цінність навчального контенту~\cite{Goncharenko2011Pedagogy}.

Принцип \emph{практичної спрямованості} орієнтує підготовку студентів на створення освітніх продуктів, придатних до безпосереднього використання в реальному освітньому середовищі, що передбачає відпрацювання повного циклу їх розробки. Дослідники акцентують увагу на автоматизації процесів, гейміфікації та прикладному характері імерсивних ресурсів як чинниках посилення практичної орієнтації навчання~\cite{Dychkivska2012Innovative}. С.~Литвинова та співавтори зазначають, що використання AR/VR-технологій уможливлює безпечне моделювання потенційно небезпечних хімічних і фізичних процесів~\cite[c.~18; 55]{Lytvynova2023Monograph}, підкреслюючи, що «за допомогою VR та AR технологій можна провести дослід з небезпечними хімічними речовинами і при цьому не завдати шкоди ні собі, ні оточенню»~\cite[c.~30]{Lytvynova2023Monograph}. У цьому контексті науковці наголошують, що практична спрямованість навчання сприяє формуванню готовності здобувачів освіти до професійної діяльності~\cite{Dychkivska2012Innovative}.

\emph{Принцип креативності} стимулює пошук нових форм моделювання навчальних ситуацій. Використання інноваційних засобів подання матеріалу зумовило якісне зрушення. Студенти генерують оригінальні ідеї в проектах. С.~Литвинова та співавтори описують процес створення власних AR-об'єктів як потужний інструмент розвитку креативності~\cite[c.~56--57]{Lytvynova2023Monograph}. У педагогічній літературі зазначається, що креативність посилюється через взаємодію з інтелектуальними системами та імерсивними середовищами~\cite{Dychkivska2012Innovative}. Дослідники обґрунтовують провідну роль інформаційно-комунікаційних технологій у розвитку креативного мислення майбутніх учителів інформатики, зокрема через застосування евристичних методів (мозковий штурм, синектика та проекти), які уможливлюють колективне генерування ідей~\cite{Mintii2023Ref}.

\emph{Принцип варіативності} відкриває можливості використання різних платформ, технологій і форматів створення імерсивного контенту з урахуванням індивідуальних потреб і освітніх запитів здобувачів. С.~Литвинова та співавтори пов'язують реалізацію цього принципу з побудовою індивідуальних освітніх траєкторій і адаптивністю імерсивних ресурсів, наголошуючи, що «індивідуальна освітня траєкторія \ldots{} ґрунтується на виборі здобувачем освіти видів, форм і темпу здобуття освіти»~\cite[c.~5; 27]{Lytvynova2023Monograph}. У ширшому педагогічному вимірі варіативність розглядається як ключова умова реалізації особистісно орієнтованого підходу, що передбачає створення освітнього середовища, здатного забезпечити свободу вибору змісту, форм і способів навчальної діяльності з метою самореалізації та професійного зростання майбутнього вчителя~\cite{Dychkivska2012Innovative}. У цьому контексті варіативність імерсивних освітніх ресурсів постає не лише технологічною, а й дидактично значущою характеристикою підготовки майбутніх педагогів.

Принцип інтеграції підкреслює необхідність узгодження ІТ-компонентів із педагогічним змістом і психологічними закономірностями навчання, що забезпечує цілісність освітнього процесу. Поєднання VR/AR із традиційними інформаційно-комунікаційними технологіями сприяє глибокій інтеграції цифрових і дидактичних елементів навчання. С.~Литвинова та співавтори визначають «інтеграцію віртуального контенту з фізичним середовищем» як одну з ключових переваг імерсивних технологій~\cite[c.~6]{Lytvynova2023Monograph}. У методичному вимірі науковці наголошують на принципі інтеграції педагогічних і цифрових технологій як умові створення електронного освітнього середовища~\cite{Lytvynova2023Monograph}. Концептуально споріднену позицію обґрунтовують П.~Мішра та М.~Коелер~\cite{Gurevych2015Integration}, які у рамках концепції TPACK (Technological Pedagogical Content Knowledge) доводять, що ефективна інтеграція технологій у навчання потребує взаємопроникнення технологічних, педагогічних і предметних знань на всіх рівнях підготовки вчителя. У методиці навчання розробки імерсивних освітніх ресурсів принцип інтеграції реалізується через поєднання педагогічного дизайну, цифрових інструментів і психолого-дидактичних вимог на всіх етапах проєктування ІОР, що забезпечує узгодженість технологічних рішень із навчальними цілями, змістом і очікуваними результатами освітньої діяльності.

Принцип \emph{рефлексивності} спрямований на розвиток здатності до усвідомленого аналізу власної діяльності та результатів проєктування імерсивних освітніх ресурсів, що забезпечує цілеспрямоване вдосконалення VR/AR-рішень. Ключову роль у цьому процесі відіграє зворотний зв'язок як основа формувального оцінювання та корекції технологічних і дидактичних рішень. С.~Литвинова та співавтори розглядають аналіз характеристик об'єктів, рефлексію й формувальне оцінювання як невід'ємні складники роботи з імерсивними ресурсами, наголошуючи на необхідності «продумати форму зворотного зв'язку з учнями»~\cite[c.~28; 37]{Lytvynova2023Monograph}.

У психолого-педагогічному вимірі рефлексивність трактується як принцип, що забезпечує здатність майбутнього педагога до самоаналізу, оцінювання досягнень і виявлення проблемних зон діяльності, стимулюючи професійне самовдосконалення~\cite{ReflexiveApproach2018}. У методичному контексті підкреслюється необхідність систематичного осмислення результатів використання віртуальної наочності та ефективності цифрових інструментів як умови формування методичних умінь і професійної готовності вчителя~\cite{Lytvynova2023Monograph}.

У межах методики навчання розробки імерсивних освітніх ресурсів принцип рефлексивності реалізується через систематичний самоаналіз проєктних рішень і оцінювання відповідності створених ІОР навчальним цілям та очікуваним результатам.

\emph{Принцип безперервності} відображає потребу в постійному оновленні цифрових і методичних компетентностей майбутнього вчителя в умовах стрімкого розвитку імерсивних технологій. У методиці навчання розробки імерсивних освітніх ресурсів цей принцип реалізується через поєднання формалізованого навчання з самонавчанням, мікрокурсами та поетапною проєктною діяльністю, що забезпечує поступове ускладнення завдань і накопичення професійного досвіду. С.~Литвинова та співавтори підкреслюють, що доступність імерсивних ресурсів у будь-якому місці й часі актуалізує парадигму «навчання впродовж життя»~\cite[c.~4; 21]{Lytvynova2023Monograph}. У свою чергу дослідники наголошують, що безперервність професійного розвитку є ключовою умовою адаптації педагога до швидких змін цифрових технологій та освітніх середовищ~\cite{Sysoieva2015Pedagogy}.

Узагальнено сукупність дидактичних підходів та принципів навчання розробки імерсивних освітніх ресурсів майбутніх учителів інформатики представлено на рис.~\ref{fig:2.2}.

\begin{figure}[!h]
\centering
\resizebox{\textwidth}{!}{%
\begin{tikzpicture}
\small

% Title
\node[dia/header=16cm] (title) {Методологічні засади навчання розробки імерсивних освітніх ресурсів};

% Column headers — centers at -6.6, -2.2, +2.2, +6.6 (4.4cm spacing, 3.0cm boxes)
\node[dia/rbox=3.0cm, minimum height=1.2cm, below=1.8cm of title.south, xshift=-6.6cm] (h1)
  {\textbf{Ціннісно-світоглядні підходи}};
\node[dia/rbox=3.0cm, minimum height=1.2cm, below=1.8cm of title.south, xshift=-2.2cm] (h2)
  {\textbf{Процесуально-діяльнісні підходи}};
\node[dia/rbox=3.0cm, minimum height=1.2cm, below=1.8cm of title.south, xshift=2.2cm] (h3)
  {\textbf{Технологічно-інноваційні підходи}};
\node[dia/rbox=3.0cm, minimum height=1.2cm, below=1.8cm of title.south, xshift=6.6cm] (h4)
  {\textbf{Системно-структурні підходи}};

% Bullet lists — all anchored to same y as l2 (tallest list)
\node[dia/lrbox=3.0cm, below=0.4cm of h2] (l2) {%
  \textbullet~діяльнісний\\
  \textbullet~конструктивістський\\
  \textbullet~контекстний\\
  \textbullet~рефлексивний%
};
\node[dia/lrbox=3.0cm, anchor=north] at (h1 |- l2.north) (l1) {%
  \textbullet~аксіологічний\\
  \textbullet~компетентнісний%
};
\node[dia/lrbox=3.0cm, anchor=north] at (h3 |- l2.north) (l3) {%
  \textbullet~інноваційний\\
  \textbullet~когнітивний\\
  \textbullet~синергетичний%
};
\node[dia/lrbox=3.0cm, anchor=north] at (h4 |- l2.north) (l4) {%
  \textbullet~системний\\
  \textbullet~інтегративний%
};

% Lines from title to column headers
\draw[dia/line] (title.south) -- ++(0,-0.5cm) -| (h1.north);
\draw[dia/line] (title.south) -- ++(0,-0.5cm) -| (h2.north);
\draw[dia/line] (title.south) -- ++(0,-0.5cm) -| (h3.north);
\draw[dia/line] (title.south) -- ++(0,-0.5cm) -| (h4.north);

% Arrows from headers to bullet lists
\draw[dia/arrow] (h1.south) -- (l1.north);
\draw[dia/arrow] (h2.south) -- (l2.north);
\draw[dia/arrow] (h3.south) -- (l3.north);
\draw[dia/arrow] (h4.south) -- (l4.north);

% Principle boxes — all at same y (below the tallest list l2)
\node[dia/rbox=3.0cm, below=0.6cm of l2] (p2a) {інтерактивність};
\node[dia/rbox=3.0cm, anchor=north] at (h1 |- p2a.north) (p1a) {науковість};
\node[dia/rbox=3.0cm, anchor=north] at (h3 |- p2a.north) (p3a) {варіативність};
\node[dia/rbox=3.0cm, anchor=north] at (h4 |- p2a.north) (p4a) {практична спрямованість};

\node[dia/rbox=3.0cm, below=0.4cm of p2a] (p2b) {рефлексивність};
\node[dia/rbox=3.0cm, anchor=north] at (h1 |- p2b.north) (p1b) {безперервність};
\node[dia/rbox=3.0cm, anchor=north] at (h3 |- p2b.north) (p3b) {креативність};
\node[dia/rbox=3.0cm, anchor=north] at (h4 |- p2b.north) (p4b) {інтеграція};

% Arrows from bullet lists to principle boxes
\draw[dia/arrow] (l1.south) -- (p1a.north);
\draw[dia/arrow] (l2.south) -- (p2a.north);
\draw[dia/arrow] (l3.south) -- (p3a.north);
\draw[dia/arrow] (l4.south) -- (p4a.north);

\end{tikzpicture}%
}

\caption{Методологічні засади дослідження}
\label{fig:2.2}
\end{figure}

Запропоновані методологічні засади охоплюють одинадцять дидактичних підходів, що відображають багатовимірний характер процесу навчання розробки імерсивних освітніх ресурсів майбутніх учителів інформатики. З метою забезпечення цілісності та внутрішньої логіки методологічної основи означені підходи згруповано у чотири кластери, обґрунтування яких спирається на фактори моделі CAMIL~\cite{Makransky2021CAMIL} та компоненти моделі TPACK~\cite{Gurevych2015Integration}.

\emph{Перший кластер}~-- \emph{ціннісно-світоглядні підходи}~-- \emph{об'єднує аксіологічний і компетентнісний підходи}, що адресують фактори \emph{мотивації} та \emph{саморегуляції} моделі CAMIL. Аксіологічний підхід задає нормативно-ціннісні рамки використання цифрових технологій і формує внутрішню мотивацію до створення якісного імерсивного контенту, тоді як компетентнісний визначає безперервну природу професійної підготовки та забезпечує саморегуляцію через усвідомлення власного рівня TPACK-компетентності. Саме з цього кластера логічно випливають принципи науковості та безперервності, які забезпечують фундаментальність і педагогічну вмотивованість розробки VR/AR-ресурсів.

\emph{Другий кластер}~-- \emph{процесуально-діяльнісні підходи}~-- \emph{включає діяльнісний, конструктивістський, контекстний і рефлексивний підходи}, які адресують фактори \emph{агентності}, \emph{втілення} та \emph{присутності} моделі CAMIL і відповідають циклу досвідного навчання Колба~\cite{Kee2023Experiential}: діяльнісний підхід забезпечує конкретний досвід (CE), конструктивістський~-- активне експериментування (AE), контекстний~-- наближення до реальних педагогічних умов через автентичну практику, а рефлексивний~-- рефлексивне спостереження (RO) та абстрактну концептуалізацію (AC). Кластер зумовлює принципи інтерактивності та рефлексивності, що відображають активну позицію здобувача у створенні імерсивного контенту.

\emph{Третій кластер}~-- \emph{технологічно-інноваційні підходи}~-- \emph{охоплює інноваційний, когнітивний і синергетичний підходи}, які адресують фактор \emph{когнітивного навантаження} моделі CAMIL та TK-компонент моделі TPACK. Інноваційний підхід визначає потребу у гнучкому включенні новітніх VR/AR-інструментів, когнітивний~-- у свідомому управлінні когнітивним навантаженням при проєктуванні імерсивних ресурсів~\cite{Buchner2021CogLoad}, синергетичний~-- у взаємодії компонентів навчального середовища, що породжують емерджентні освітні ефекти. Логічними принципами цього кластера є варіативність і креативність, які забезпечують здатність майбутнього вчителя працювати з різними технологічними платформами та створювати оригінальні імерсивні рішення.

\emph{Четвертий кластер}~-- \emph{системно-структурні підходи}~-- \emph{об'єднує системний та інтегративний підходи}, що адресують \emph{TPK-інтеграцію} моделі TPACK та забезпечують системну узгодженість шестикомпонентної моделі проєктування ІОР (підрозділ~1.2). Системний підхід забезпечує узгодженість усіх елементів освітнього процесу, а інтегративний~-- поєднання технологічних, дидактичних і психологічних компонентів у єдиний зміст підготовки~\cite{Schmidt2024EntangledLXD}. Цей кластер формує основу для принципів практичної спрямованості та інтегративності, що характеризують підготовку як процес цілісного розвитку професійних компетентностей вчителя.

Згруповані підходи утворюють логічно вмотивовану структуру, з якої випливають відповідні принципи навчання. Теоретична обґрунтованість кластеризації підтверджується матрицею відповідності (рис.~\ref{fig:approaches_matrix}): кожен кластер забезпечує покриття відповідних факторів CAMIL, а сукупність чотирьох кластерів адресує всі шість медіативних факторів та усуває виявлений у підрозділі~1.1 дефіцит TPK-компетентності.

Запропонована методологічна основа знаходить безпосереднє втілення у практичній роботі автора. Навчальний посібник з WebAR-розробки~\cite{Shepiliev2020WebAR,Shepiliev2021Career} реалізує діяльнісний і конструктивістський підходи через систему практичних завдань зі створення імерсивних ресурсів; методика інтеграції моделей машинного навчання у WebAR~\cite{Semerikov2024WebARML_CoSinE} втілює інноваційний і когнітивний підходи; структура навчального курсу «Синтетична реальність в освіті» відображає системний та інтегративний підходи через поетапну організацію від базових технологій до комплексних проєктів.

Узагальнення дидактичних підходів та їх кластеризація дали змогу структурувати методологічні засади навчання розробки імерсивних освітніх ресурсів. Виокремлені чотири кластери підходів визначають ціннісні, діяльнісні, технологічні та системно-структурні виміри підготовки майбутніх учителів інформатики, з яких логічно випливають принципи навчання. Сукупно означені методологічні орієнтири формують концептуальне підґрунтя для подальшого конструювання моделі формування вмінь створювати імерсивні освітні ресурси у процесі фахової підготовки, розроблення якої представлено в наступному підрозділі.

